\documentclass[11pt,openany,a4paper]{memoir}
\usepackage[T1]{fontenc}
\usepackage[utf8]{inputenc}
\usepackage{mathpazo}
\usepackage{microtype}
\newcommand{\olpath}{../upstream}
\input{\olpath/sty/open-logic.sty}
% Javanese text layer; Babel's English backend is used only for package compatibility.
% All displayed captions and token forms used in this tranche are overridden.
\setlocalecaption{english}{theorem}{Teorema}
\setlocalecaption{english}{example}{Tuladha}
\setlocalecaption{english}{definition}{Definisi}
\setlocalecaption{english}{lemma}{Lema}
\setlocalecaption{english}{proposition}{Proposisi}
\setlocalecaption{english}{corollary}{Korolari}
\setlocalecaption{english}{problem}{Soal}
\setlocalecaption{english}{problems}{Soal}
\setlocalecaption{english}{remark}{Cathetan}
\setlocalecaption{english}{note}{Cathetan}
\setlocalecaption{english}{axiom}{Aksioma}
\setlocalecaption{english}{case}{Kasus}
\setlocalecaption{english}{convention}{Konvensi}
\settexttoken{element}{anggota}{anggota}[Anggota][Anggota]
\definetoken{a}{element}{salah siji}
\definetoken{A}{element}{Salah siji}
\settexttoken{formula}{formula}{formula}[Formula][Formula]
\settexttoken{derivation}{derivasi}{derivasi}[Derivasi][Derivasi]
% The source's identity token names the predicate, rather than equality in general.
\settexttoken{identity}{predikat identitas}{predikat identitas}[Predikat identitas][Predikat identitas]
\settexttoken{conditional}{kondisional}{kondisional}[Kondisional][Kondisional]
\settexttoken{biconditional}{bikondisional}{bikondisional}[Bikondisional][Bikondisional]
\settexttoken{falsity}{kasalahan}{kasalahan}[Kasalahan][Kasalahan]
\settexttoken{truth}{kabeneran}{kabeneran}[Kabeneran][Kabeneran]
\settexttoken{derivability}{derivabilitas}{derivabilitas}[Derivabilitas][Derivabilitas]
% The source token layer distinguishes logical symbols from their interpretations.
\settexttoken{language}{basa}{basa}[Basa][Basa]
\settexttoken{subformula}{subformula}{subformula}[Subformula][Subformula]
\settexttoken{sentence}{ukara}{ukara}[Ukara][Ukara]
\settexttoken{variable}{variabel}{variabel}[Variabel][Variabel]
\settexttoken{propositional variable}{variabel proposisional}{variabel proposisional}[Variabel proposisional][Variabel proposisional]
\settexttoken{constant}{simbol konstanta}{simbol konstanta}[Simbol konstanta][Simbol konstanta]
\settexttoken{predicate}{simbol predikat}{simbol predikat}[Simbol predikat][Simbol predikat]
\settexttoken{function}{simbol fungsi}{simbol fungsi}[Simbol fungsi][Simbol fungsi]
\settexttoken{operator}{operator logis}{operator logis}[Operator logis][Operator logis]
\settexttoken{main operator}{operator utama}{operator utama}[Operator utama][Operator utama]
\settexttoken{structure}{struktur}{struktur}[Struktur][Struktur]
\settexttoken{valuation}{valuasi}{valuasi}[Valuasi][Valuasi]
\settexttoken{domain}{ranah}{ranah}[Ranah][Ranah]
\settexttoken{value}{nilai}{nilai}[Nilai][Nilai]
\settexttoken{tableau}{tableau}{tableau}[Tableau][Tableau]
\settexttoken{signed formula}{formula mawa tandha}{formula mawa tandha}[Formula mawa tandha][Formula mawa tandha]
\settexttoken{free for}{bebas kanggo}{bebas kanggo}[Bebas kanggo][Bebas kanggo]
\settexttoken{derive}{nderivasi}{nderivasi}[Nderivasi][Nderivasi]
\settexttoken{discharge}{mbubarake}{mbubarake}[Mbubarake][Mbubarake]
\settexttoken{discharged}{dibubarake}{dibubarake}[Dibubarake][Dibubarake]
\settexttoken{undischarged}{durung dibubarake}{durung dibubarake}[Durung dibubarake][Durung dibubarake]
\settexttoken{nonderivability}{kahanan ora bisa diderivasi}{kahanan ora bisa diderivasi}[Kahanan ora bisa diderivasi][Kahanan ora bisa diderivasi]
\definetoken{a}{language}{salah siji}
\definetoken{A}{language}{Salah siji}
\definetoken{a}{subformula}{salah siji}
\definetoken{A}{subformula}{Salah siji}
\definetoken{a}{sentence}{salah siji}
\definetoken{A}{sentence}{Salah siji}
\definetoken{a}{variable}{salah siji}
\definetoken{A}{variable}{Salah siji}
\definetoken{a}{propositional variable}{salah siji}
\definetoken{A}{propositional variable}{Salah siji}
\definetoken{a}{constant}{salah siji}
\definetoken{A}{constant}{Salah siji}
\definetoken{a}{predicate}{salah siji}
\definetoken{A}{predicate}{Salah siji}
\definetoken{a}{function}{salah siji}
\definetoken{A}{function}{Salah siji}
\definetoken{a}{operator}{salah siji}
\definetoken{A}{operator}{Salah siji}
\definetoken{a}{main operator}{salah siji}
\definetoken{A}{main operator}{Salah siji}
\definetoken{a}{structure}{salah siji}
\definetoken{A}{structure}{Salah siji}
\definetoken{a}{valuation}{salah siji}
\definetoken{A}{valuation}{Salah siji}
\definetoken{a}{domain}{salah siji}
\definetoken{A}{domain}{Salah siji}
\definetoken{a}{value}{salah siji}
\definetoken{A}{value}{Salah siji}
\definetoken{a}{tableau}{salah siji}
\definetoken{A}{tableau}{Salah siji}
\definetoken{a}{signed formula}{salah siji}
\definetoken{A}{signed formula}{Salah siji}
\definetoken{a}{formula}{salah siji}
\definetoken{A}{formula}{Salah siji}
\definetoken{a}{derivation}{salah siji}
\definetoken{A}{derivation}{Salah siji}
\settexttoken{surjective}{surjektif}{surjektif}[Surjektif][Surjektif]
\settexttoken{injective}{injektif}{injektif}[Injektif][Injektif]
\settexttoken{bijective}{bijektif}{bijektif}[Bijektif][Bijektif]
\definetoken{a}{surjective}{}
\definetoken{A}{surjective}{}
\definetoken{a}{injective}{}
\definetoken{A}{injective}{}
\definetoken{a}{bijective}{}
\definetoken{A}{bijective}{}
\settexttoken{surjection}{surjeksi}{surjeksi}[Surjeksi][Surjeksi]
\settexttoken{injection}{injeksi}{injeksi}[Injeksi][Injeksi]
\settexttoken{bijection}{bijeksi}{bijeksi}[Bijeksi][Bijeksi]
\definetoken{a}{surjection}{salah siji}
\definetoken{A}{surjection}{Salah siji}
\definetoken{a}{injection}{salah siji}
\definetoken{A}{injection}{Salah siji}
\definetoken{a}{bijection}{salah siji}
\definetoken{A}{bijection}{Salah siji}
\addto\captionsenglish{%
  \renewcommand{\contentsname}{Isine}
  \renewcommand{\chaptername}{Bab}
  \renewcommand{\partname}{Perangan}
  \renewcommand{\figurename}{Gambar}
  \renewcommand{\tablename}{Tabel}
  \renewcommand{\proofname}{Bukti}
  \renewcommand{\bibname}{Kapustakan}
}
\addto\extrasenglish{%
  \Crefname{thm}{Teorema}{Teorema}
  \Crefname{ex}{Tuladha}{Tuladha}
  \Crefname{defn}{Definisi}{Definisi}
  \Crefname{lem}{Lema}{Lema}
  \Crefname{prop}{Proposisi}{Proposisi}
  \Crefname{prob}{Soal}{Soal}
  \Crefname{rem}{Cathetan}{Cathetan}
  \Crefname{cor}{Korolari}{Korolari}
  \Crefname{axiom}{Aksioma}{Aksioma}
  \Crefname{note}{Cathetan}{Cathetan}
  \Crefname{case}{Kasus}{Kasus}
  \Crefname{conv}{Konvensi}{Konvensi}
  \Crefname{figure}{Gambar}{Gambar}
  \Crefname{table}{Tabel}{Tabel}
  \Crefname{section}{Perangan}{Perangan}
  \Crefname{chapter}{Bab}{Bab}
  \Crefname{page}{kaca}{kaca}
}
\newcommand{\JvReferenceName}[2]{\crefname{#1}{#2}{#2}\Crefname{#1}{#2}{#2}}
\AtBeginDocument{%
  \JvReferenceName{thm}{Teorema}
  \JvReferenceName{ex}{Tuladha}
  \JvReferenceName{defn}{Definisi}
  \JvReferenceName{lem}{Lema}
  \JvReferenceName{prop}{Proposisi}
  \JvReferenceName{prob}{Soal}
  \JvReferenceName{rem}{Cathetan}
  \JvReferenceName{cor}{Korolari}
  \JvReferenceName{axiom}{Aksioma}
  \JvReferenceName{note}{Cathetan}
  \JvReferenceName{case}{Kasus}
  \JvReferenceName{conv}{Konvensi}
  \JvReferenceName{figure}{Gambar}
  \JvReferenceName{table}{Tabel}
  \JvReferenceName{section}{Perangan}
  \JvReferenceName{chapter}{Bab}
  \JvReferenceName{page}{kaca}
  \renewcommand{\crefpairconjunction}{ lan }
  \renewcommand{\crefmiddleconjunction}{, }
  \renewcommand{\creflastconjunction}{, lan }
  \renewcommand{\crefrangeconjunction}{ nganti }
}

\setlrmarginsandblock{28mm}{28mm}{*}
\setulmarginsandblock{25mm}{28mm}{*}
\checkandfixthelayout
\linespread{1.05}
\hyphenpenalty=10000
\exhyphenpenalty=10000
\tolerance=2000
\emergencystretch=5em
\pdfinfoomitdate=1
\pdftrailerid{}
\pdfsuppressptexinfo=15
\hypersetup{pdftitle={OpenLogic: Dhasar Himpunan, Logika, lan Teori Model - Basa Jawa},pdfauthor={Open Logic Project; Codex translation},pdflang={jv-Latn-ID}}
\begin{document}
\begin{titlingpage}
\raggedright
{\sffamily\Large OpenLogic\par}
\vspace{18mm}
{\sffamily\bfseries\Huge Dhasar Himpunan, Logika, lan Teori Model\par}
\vspace{6mm}
{\Large Basa Jawa, aksara Latin\par}
\vspace{18mm}
{\large Wacan kumulatif nganti bab interpolasi.\par}
\vfill
Terjemahan mesin adhedhasar Open Logic Project.
Potret sumber kanggo wacan iki: 202 saka 722 berkas sumber wis diterjemahake.
Edhisi jangkep isih digarap.
\par\medskip
Sumber Inggris: OpenLogicProject/OpenLogic, revisi
\texttt{9620cc73f9c8e0ad003c514a5d3748f29611c4c0}.
\par\medskip
Karya asli lan terjemahan iki nganggo
\href{https://creativecommons.org/licenses/by/4.0/}{Creative Commons Attribution 4.0}.
Hak lan atribusi komponen asli tetep ditrapake.
Terjemahan iki ora nuduhake panyengkuyung saka Open Logic Project.
\end{titlingpage}
\pagestyle{plain}
\chapter*{Babagan Open Logic Project}
\addcontentsline{toc}{chapter}{Babagan Open Logic Project}


\textit{Open Logic Text} yaiku buku ajar sumber kabuka kang disusun
bebarengan, ngenani metalogika formal lan metodhe formal, wiwit tataran
madya (yaiku, sawise kuliah pambuka logika formal). Sanajan ditujokake
marang pamaca kang dhasar pasinaone dudu matematika (mligine mahasiswa
filsafat lan ilmu komputer), andharane tetep tliti.

Sawetara topik kang saiki wis dirembug bisa uga durung jangkep
andharane, lan akeh perangan isih mbutuhake revisi gedhe. Ing tembe,
kita ngrancang ngrembakakake teks iki supaya ngrembug luwih akeh topik.
Kita uga ngrancang nambahake glosarium, dhaptar wacan luwih lanjut,
cathetan sejarah, gambar, andharan kang luwih cetha, perangan kang
njlentrehake gayutane asil karo filsafat, ilmu komputer, lan matematika,
sarta luwih akeh soal lan tuladha ing teks iki.
Yen nemokake kaluputan utawa duwe usul,
\href{https://github.com/OpenLogicProject/OpenLogic/wiki/Contributing}{tulung ngandhani tim proyek}.

Proyek iki lumaku kanthi semangat sumber kabuka. Teks iki ora mung bisa
diakses kanthi bebas, nanging kita uga nyedhiyakake sumber LaTeX kanthi
Lisènsi Creative Commons Attribution. Lisènsi iki menehi hak marang
sapa wae kanggo ngundhuh, nggunakake, ngowahi, nata maneh, ngowahi
format, lan ngedum maneh karya kita, anggere menehi krédhit kang patut.
Katrangan liyane bisa diwaca ing situs Open Logic Project:
\href{http://openlogicproject.org/}{openlogicproject.org}.

\clearpage
\tableofcontents*

\olimport[../translation/content/sets-functions-relations]{sets-functions-relations-complete}

% The source's propositional part also imports the first-order proof chapters.
% This reader presents those chapters once in the first-order part below.
\olpart{pl}{Logika Proposisional}
\begin{editorial}
Perangan iki ngemot bab sintaksis lan semantik logika proposisional.
Sistem bukti kang digunakake bebarengan karo logika sepisanan
kapacak sapisan ing perangan logika sepisanan ing ngisor iki.
\end{editorial}
\olimport[../translation/content/propositional-logic/syntax-and-semantics]{syntax-and-semantics}
\OLEndPartHook

\olimport[../translation/content/first-order-logic]{first-order-logic}

% Lindstrom belongs to the next, incomplete chapter and is not in this reader.
\olpart{mod}{Teori Model}
\begin{editorial}
Perangan iki ngemot bab dhasar teori model, model aritmetika,
lan interpolasi. Bab Lindstr\"om isih digarap lan durung
dilebokake ing wacan iki.
\end{editorial}
\olimport[../translation/content/model-theory/basics]{basics}
\olimport[../translation/content/model-theory/models-of-arithmetic]{models-of-arithmetic}
\olimport[../translation/content/model-theory/interpolation]{interpolation}
\OLEndPartHook

\chapter*{Cathetan Sumber lan Notasi}
\addcontentsline{toc}{chapter}{Cathetan Sumber lan Notasi}

Ing tuladha relasi tatanan, sumber nganggo $I$ tanpa
ngenali jenenge luwih dhisik. Ing kono, $I$ kudu diwaca
minangka relasi identitas ing wilangan natural,
yaiku pasangan-pasangan ing diagonal kang wis diterangake.

Ing definisi pang wit, sumber nulis $z \in X \setminus B$,
nanging himpunan simpul wit wis dijenengi $A$ lan $X$
durung ditemtokake. Ing definisi iku, wacanen $X$ minangka $A$.
Maksimalitas banjur ateges saben simpul ing njaba $B$
ora bisa dibandhingake karo paling ora siji simpul ing $B$.

Notasi $R^+$ nduweni teges lokal kang beda:
ing perangan tatanan, tegese tutupan refleksif;
ing perangan operasi relasi, tegese tutupan transitif.
Gunakna definisi ing perangan kang lagi diwaca.

Andharan ekstensionalitas ing wiwitan kudu diwaca
bebarengan karo perangan Paradoks Russell:
ekstensionalitas njamin mung ana siji himpunan
kang cocog yen himpunan kasebut pancen ana.

Ing tuladha fungsi oyod kuadrat, sumber nyebut oyod kang dipilih
minangka positif, nanging domain wilangan natural kalebu nol.
Ing nol, wacanen iki minangka oyod kuadrat utama kang ora negatif.
Iki temuan audit sumber OLFUN-002 lan wis dibenerake ing teks terjemahan.
Ing tuladha himpunan pandhisik lan panerus, masukan dijenengi $n$,
nanging banjur ditulis $x$. Ing kono, $x$ nuduhake masukan $n$ kang padha.
Iki temuan OLFUN-003; teks terjemahan saiki nggunakake $x$ kanthi ajeg.

Proposisi yen saben injeksi duwe invers kiwa mbutuhake katrangan
tumrap domain kosong. Bukti ing sumber milih $a \in A$, mula
lumaku yen $A$ ora kosong. Yen $A$ lan $B$ loro-lorone kosong,
fungsi kosong dadi inverse dhewe. Nanging yen $A$ kosong lan $B$
ora kosong, injeksi kosong ora duwe invers kiwa saka $B$ menyang $A$.
Mula, wacanen proposisi iku kanthi sarat $A$ ora kosong utawa $B$ kosong.
Iki temuan OLFUN-001; awak proposisi terjemahan nganggo sarat prasaja
$A$ ora kosong lan cathetan jejer menehi sarat umum kang pas.
Panganggone Aksioma Pilihan ing bukti invers tengen wis diterangake
ing cathetan sikil sumber.

Ing andharan graf fungsi, sumber nyebut $R_f\subseteq A\times B$
minangka relasi "ing $A\times B$". Miturut definisi buku iki,
relasi ing himpunan kasebut malah dadi subhimpunan
$(A\times B)^2$. Teks terjemahan wis mbenerake iki dadi relasi
antarane $A$ lan $B$ kang kaemot ing $A\times B$ (OLFUN-004).

Watesan fungsi ing $C$ mung nyaring koordinat masukan, mula grafe
$R_f\cap(C\times B)$. Watesan relasi ing bagean sadurunge yaiku
$R\cap C^2$ lan nyaring loro koordinat. Rumus fungsi ing sumber bener;
pratelan yen loro gagasan persis padha wis diwatesi dadi analogi
kanthi bedane dijlentrehake (OLFUN-005).

Audit sumber OLSIZ-20260904 dikunci dening REVIEW.md kanthi SHA-256
26913176baacbda5ae8a47bcc82ccf0366b0763313e6eeeb45df59e64249a1f9
lan FINDINGS.json kanthi SHA-256
9b6e836c8432eb75d331913983603796d6557da6ec3ad71b846b2a248374cd07.
Sumber Inggris tetep persis kaya asline; koreksi ing ngisor iki ditrapake
ing awak terjemahan lan diwenehi ID audit ing jejer panggonane.
Temuan tambahan kang biyen diusulake minangka OLSIZ-011 banjur ditarik
manèh sawise pamriksan bita nuduhake saben tabel wis duwe telung garis
miring walik: pamisah larik banjur prentah $\backslash$hline kang bener.
Mula ora ana koreksi tabel adhedhasar usulan kang wis ditarik iku.

Ing OLP-0029, larik nilai tabel wilangan wutuh ora ngemot $-3$
ing sangisore $f(7)$, senajan rumuse mbutuhake nilai iku. Larik terjemahan
wis nambahake $-3$ (OLSIZ-001). Ing OLP-0031, definisi kowinates saiki
nyebut komplemen ing $\Nat$ saka subhimpunan winates ing kono,
manut rumus $\Nat\setminus A$ (OLSIZ-002).

Ing OLP-0032, urutan kulawarga pasangan saiki nulis
$\tuple{2,m}$ banjur $\tuple{3,m}$, manut tabel sabanjure (OLSIZ-003).
Ing OLP-0034, rerangken karakteristik nggunakake jeneng metuane $s$
kanthi ajeg (OLSIZ-004). Tuladha $h(n)$ saiki menehi $n$ nol banjur
buntut siji tanpa wates, mula saben metuane pancen ana ing $\Bin^\omega$
(OLSIZ-005).

Ing OLP-0035, rong cabang himpunan kosong saiki nganggo $f(x)=y$,
yaiku bijeksi kang lagi dirembug, dudu enumerasi $g$ kang durung
dikenalake (OLSIZ-006). Ing OLP-0036, dudutan bukti diagonal saiki
dikuwantifikasi kanggo saben $x\in A$, saengga kabeh jangkauan $g$
kacetha ora ngemot himpunan diagonal (OLSIZ-007).

Ing OLP-0039, $s_n(m)$ diterangake minangka digit kaping $m$ saka
rerangken kaping $n$ (OLSIZ-008), lan pandhuan diagonal ngganti $1$
dadi $0$ sarta $0$ dadi $1$ (OLSIZ-009). Limang prentah garis horisontal
asli tetep dijaga amarga bita-bita TeX-e wis bener.
Ing OLP-0040, rerangken karakteristik uga nggunakake jeneng $s$
tanpa indeks kang ora kaiket (OLSIZ-010).

Ing OLP-0043, andharan tatanan luwih dhisik kanthi bener nggunakake
$s-r$ kanggo nemtokake $r\leq s$, nanging ukara ``yaiku'' sabanjure
malah mbalikke bedane dadi $r-s$. Rumus kang ditampilake nggunakake
$s-r$ maneh. Teks terjemahan manut urutan kang ajeg iki lan menehi
cathetan koreksi ing jejer rumus.

Ing OLP-0044, ana tandha kurang-saka kang kesasar sanalika sawise
token jamak kanggo unsur ing ukara bab himpunan rasional ing sangisore
$\sqrt{2}$. Tandha iku ora duwe guna sintaktis utawa matematis lan
ora digawa menyang teks terjemahan; koreksine dicathet ing jejer ukara.
Cathetan sikil kang dipateni ing sumber uga kandha yen simbol
$\sqrt{2}$ nuduhake oyod positif lan negatif kanthi padha. Miturut
konvensi baku kang digunakake teorema, simbol iku nuduhake oyod utama
kang ora negatif. Komentar iku tetep dipateni lan diterjemahake mung
kanggo njaga kaselarasan sumber.

Ing OLP-0045, butir kapisan bukti kajangkepan kandha yen anane wates
ndhuwur njamin paling ora siji potongan kalebu ing $S$. Anane potongan
iku sejatine dumadi amarga $S$ wis dianggep ora kosong. Sarat wates
ndhuwur lagi digunakake sabanjure kanggo nuduhake yen gabungane dadi
subhimpunan murni saka wilangan rasional. Teks terjemahan nggunakake
premis kang bener lan menehi cathetan ing jejer butir bukti.

Ing OLP-0047, ukara bab wilangan real ngilangi tembung aran sawise
frasa Inggris ``the (tedious) of checking''. Terjemahan mulihake tugas
kang dimaksud, yaiku mriksa yen wilangan real mbentuk lapangan katata,
lan nyathet koreksi ing jejer ukara. Perangan iku uga njaluk pamaca
mriksa yen pangurangan lan pambagian potongan ngasilake potongan,
nanging rumus loro operasi kasebut ing OLP-0045 mung ana minangka
komentar kang dipateni. Terjemahan njaga panjaluk kasebut lan nerangake
yen panjaluk iku nganggep rumus kang dimaksud mau.

OLP-0048 duwe sawetara ukara Inggris kang kurang utawa mbaleni tembung:
``can be equally be encoded'', ``a simple way to see this as follows'',
``define we a notion'', lan ``hone on in''. Terjemahan nggunakake wacan
gramatikal kang paling cilik tanpa ngganti pratelan matematise. Sumber
uga salah nyebut wilangan real minangka relasi ekuivalensi, kamangka
ukara sabanjure kanthi bener nggunakake kelas ekuivalensi. Pratelan
lapangan katata lan kajangkepan nganggo ``rerangken Cauchy'' minangka
cekakan kanggo kelas-kelas ekuivalensine; teks Jawa nerangake wacan
quotient iki ing jejer teorema.

Definisi positif ing OLP-0048 mbandhingake kelas real karo nol rasional,
nanging ukara sabanjure nggunakake nol real lan obyek kang dibangun
pancen kalebu quotient real. Terjemahan nggunakake nol real ing loro
panggonan. Gambaran bukti kajangkepan uga nganggep saben kulawarga real
kang diwatesi duwe wates rasional konstan saka ndhuwur lan ngisor.
Iki kasunyatan Archimedes kang durung dibuktekake ing perangan kasebut;
ketergantungan iki dicathet sadurunge argumen pambagean loro diterusake.

Cathetan iki tambahan panyunting ing terjemahan mesin.
Rumus ing sumber selaras tetep dijaga supaya
bisa dibandhingake langsung karo sumber Inggris.

\bibliographystyle{natbib-oup}
\bibliography{open-logic}
\end{document}
